\documentclass[12pt]{article}
\usepackage[utf8]{inputenc}
\usepackage[T1]{fontenc}
\usepackage{hyperref}
\title{Santiago Xanica Zapotec Language Resources for K-12 Curriculum Development}
\author{Compiled by Claude Sonnet 5.5 (claude-sonnet-5-5) for Andr\'as Kornai}
\date{}
\begin{document}
\maketitle

\section{Introduction}
Santiago Xanica Zapotec (ISO 639-3 zpr; Glottocode sant1452) is an Otomanguean language of Oaxaca, Mexico. Wikipedia cites about 2,500 speakers (1990 census) and Joshua Project lists 4,000. The search found only catalogue pages, scripture portions and evangelistic audio. It found no grammar, dictionary or school materials.

\section{Primary Educational Resources}
No primers, textbooks or school programmes were found.

\section{Linguistic and Reference Materials}
\begin{itemize}
\item Wikipedia article on X\'anica Zapotec: \url{https://en.wikipedia.org/wiki/X%C3%A1nica_Zapotec}
\item Glottolog entry: \url{https://glottolog.org/resource/languoid/id/sant1452.bigmap.html}
\item The Archive of the Indigenous Languages of Latin America (AILLA) reportedly holds a ``Zapotec Survey from Santiago Xanica''. The search showed no direct link and no details of its contents.
\end{itemize}

\section{Religious and Community Resources}
\begin{itemize}
\item Joshua Project profile, which reports Bible portions from 1996--1998 but no New Testament: \url{https://joshuaproject.net/index.php/languages/zpr}
\end{itemize}

\section{Digital and Audio Resources}
\begin{itemize}
\item Global Recordings Network audio Bible stories and evangelistic messages, with MP3 and slideshow downloads: \url{https://www.globalrecordings.net/en/language/zpr}
\item Wikitongues page, which currently lists no videos: \url{https://wikitongues.org/languages/zpr/}
\end{itemize}

\section{Curriculum Development Context}
The language is endangered, and not all young people speak it. It is not known to be taught in schools. Documentation is thin, so any curriculum would start with community-led documentation and orthography work.
\end{document}